\section{Controlled HTTP executions}\label{sec:http}
\code{scripts/run_stack_experiments.py} runs six configurations, with the
versions retained in each trace: Express 5.3.0 (\code{etag} 1.8.1) with weak
and strong ETags; Python's \code{http.server} validating through
\code{Last-Modified}, with advancing and restored modification times; and
Apache with \code{FileETag MTime Size}, with advancing and restored
modification times. Every server binds only loopback. The two pages move
through the phases (eligible, 100), (ineligible, 10) and (eligible, 100); the
target pair (eligible, 10) is read by capturing admission in the first phase
and fee twice in the middle phase, then admission again. All four reads are
full GETs whose bodies are compared byte for byte with the targets; only then
does a fifth request send admission's original \code{If-None-Match}, or
\code{If-Modified-Since} without an ETag. A bodyless 304 never substitutes for
reading the endpoint body.

The two files do not update atomically as a pair, so the order of the
single-file updates matters. The forward phase replaces admission before fee,
withdrawing eligibility before introducing the lower fee; restoration replaces
fee before admission. Each update writes a staged file and renames it, and both
pages are read back after every replacement; the trace asserts from these
readbacks that the target pair was never served. Restoring admission first
would briefly serve the target pair, and the run would no longer show
noncoexistence. Fresh modification times advance by 10\,s per phase; in the
restored cases the final time equals the original, an explicit mechanism
condition rather than a claim about deployments.

\begin{table}[t]
\caption{The six HTTP configurations (macOS) and the Linux re-execution.}
\label{tab:http}
\small
\setlength{\tabcolsep}{4pt}
\begin{tabular}{@{}lllll@{}}
\toprule
Configuration & Repeats & Status & Pair served & Linux \\
\midrule
Express, weak ETag & yes & 304 & no & same \\
Express, strong ETag & yes & 304 & no & same \\
Python, advancing mtime & no & 200 & no & same \\
Python, restored mtime & yes & 304 & no & same \\
Apache, advancing mtime & no & 200 & no & same \\
Apache, restored mtime & yes & 304 & no & same \\
\bottomrule
\end{tabular}

\end{table}

In four configurations the endpoint validator repeats and the conditional GET
returns 304 although the target pair never coexisted (Table~\ref{tab:http}).
These outcomes accuse no server of violating HTTP: a 304 says that a stored
copy may be reused~\cite{http}. They test the consumer's inference from equal
validators to continuity, and 4 of 6 is not a failure rate. On 2026-11-05 the
six configurations were re-executed on Linux x86\_64 (Express 5.3.0, Python
3.13.18, Ubuntu Apache 2.4.58) with identical outcomes; the runner refuses to
overwrite traces and aborts on a failed start or a violated state assertion.
